\title{The Era of 1-bit LLMs: All Large Language Models are in 1.58 Bits}

\begin{abstract}
Emerging advancements like BitNet~\cite{bitnet} are ushering in a transformative period characterized by 1-bit Large Language Models (LLMs). In this study, we propose a 1-bit LLM version termed \textbf{\text{BitNet b1.58}}, wherein each individual parameter (or weight) of the model assumes a value from the ternary set \{-1, 0, 1\}. This model achieves parity with full-precision Transformer LLMs (utilizing FP16 or BF16) of equivalent architecture and training data volume, excelling in both perplexity and downstream performance metrics. Notably, BitNet b1.58 demonstrates significant improvements in computational efficiency, exhibiting reduced latency, memory usage, throughput demands, and energy overhead. Beyond immediate efficiency gains, the introduction of the 1.58-bit LLM establishes a novel scaling framework and methodology for developing future LLMs that are both performant and resource-conscious. Additionally, this advancement lays the foundation for a new computational approach and stimulates interest in hardware architectures purpose-built for 1-bit LLMs.
\end{abstract}

\section{The Era of 1-bit LLMs}

The past few years have witnessed exponential expansion in both the scale and capability of Large Language Models (LLMs) within artificial intelligence research. While these systems deliver outstanding results across a spectrum of natural language processing benchmarks, their ever-increasing parameter counts present significant barriers for real-world deployment and intensify worries regarding their substantial environmental and financial resource requirements, especially from elevated energy usage.

To mitigate these obstacles, post-training quantization techniques have emerged as a promising solution, facilitating the development of low-bitwidth models for inference~\cite{smoothquant, gptq, quip, quip_sharp}. By lowering the precision of both model weights and activations, these approaches offer considerable savings in memory and computational expenditure for LLMs. The field has observed a progressive shift from standard 16-bit precision toward more compact forms, notably down to 4-bit representations~\cite{gptq, awq}. Nonetheless, while post-training quantization is prevalent in industrial applications of LLMs, it often results in sub-optimal model performance compared to alternative training strategies.

Recent advancements in 1-bit model architectures, exemplified by BitNet~\cite{bitnet}, represent a compelling approach for decreasing the operational costs of LLMs without sacrificing their accuracy. Standard LLMs employ 16-bit floating-point representations (such as FP16 or BF16), and matrix multiplication operations dominate their computation workload. The predominant sources of computation overhead in these models stem from floating-point multiplications and additions. By contrast, BitNet confines its matrix multiplications to integer additions, which yields substantial energy reductions. Since the maximal computational throughput of hardware accelerators is often constrained by power limits, lowering energy requirements also enables improved computational speeds.

Beyond computational overhead, another key expense is incurred when loading model parameters from DRAM to the on-chip memory (such as SRAM) during inference. Although scaling up SRAM capacity has been proposed to enhance throughput, such modifications tend to escalate costs significantly relative to DRAM. When compared with full-precision models, 1-bit LLMs are advantageous in terms of memory efficiency, benefiting from reduced requirements in both memory size and bandwidth. These improvements substantially lower the expenses and delays associated with parameter loading, accelerating inference.

Here, we introduce a notable 1-bit LLM variant, termed \textbf{\text{BitNet b1.58}}, in which all parameters assume ternary values from the set \{-1, 0, 1\}. We extend the initial BitNet’s binary regime by incorporating an explicit zero, resulting in a 1.58-bit model in binary terms. \text{BitNet b1.58} maintains all advantages characteristic of the original 1-bit BitNet, including its efficient computation paradigm that eschews most multiplication operations in favor of optimizable approaches. Energy demands remain on par with those of BitNet, while memory footprint, throughput, and latency surpass FP16 LLM baselines in efficiency. Moreover, \text{BitNet b1.58} introduces two further strengths: first, its modeling expressiveness improves via feature selection enabled by the zero parameter, leading to significant performance boosts in the context of 1-bit LLMs; second, our empirical findings indicate that, at the 3B model scale, \text{BitNet b1.58} matches the perplexity and downstream performance of FP16 models under equivalent configuration settings (such as parameter count and training data volume).

\section{BitNet b1.58}

\text{BitNet b1.58} is derived from the BitNet model, which modifies the standard Transformer architecture by substituting \emph{nn.Linear} layers with \emph{BitLinear} layers. This model is trained from the beginning using 1.58-bit precision for weights and 8-bit precision for activations. While maintaining the foundation of BitNet, it incorporates several updates described below.

\paragraph{Quantization Function.} To ensure the weights are restricted to -1, 0, or +1, we implement the \emph{absmean} quantization approach. In this method, the weight matrix is first normalized by dividing by the mean of its absolute values, after which each element is rounded to the closest value in the set \{-1, 0, +1\}:

\begin{equation}
    \widetilde{W} = \mathrm{RoundClip}(\frac{W}{\gamma + \epsilon}, -1, 1),
\end{equation}

\begin{equation}
    \mathrm{RoundClip}(x, a, b) = \max(a, \min(b, \mathrm{round}(x))),
\end{equation}

\begin{equation}
    \gamma = \frac{1}{nm}\sum_{ij} |W_{ij}|.
\end{equation}

For activations, we adhere to the quantization method utilized in BitNet, with one significant change: activations are not normalized to the interval $[0, Q_b]$ before applying non-linear operations. Rather, for each token, activations are scaled to $[-Q_b, Q_b]$, effectively eliminating the need for zero-point quantization. This alteration streamlines both the implementation and system-level optimization, while our empirical results indicate it has minimal impact on model performance.

\paragraph{LLaMA-alike Components.}

The LLaMA architecture~\cite{llama, llama2} has established itself as a standard framework for open-source large language models. To facilitate alignment with the open-source community, we design \text{BitNet b1.58} to be compatible with LLaMA-like components. This involves utilizing RMSNorm~\cite{rmsnorm}, SwiGLU~\cite{swiglu}, rotary embeddings~\cite{rope}, and removing all bias terms. With these choices, \text{BitNet b1.58} is readily deployable within widely used open-source tools (e.g., Huggingface, vLLM~\cite{vllm}, and llama.cpp\footnote{\url{https://github.com/ggerganov/llama.cpp}}) with minimal modification required.

\section{Results}

We conducted a comparison between \text{BitNet b1.58} and our FP16 implementation of \text{LLaMA LLM} at several parameter scales. Both models were pre-trained on the RedPajama dataset~\cite{redpajama}, each processed for a total of 100 billion tokens to maintain parity in pre-training. For evaluation, we assessed the models in zero-shot settings across multiple language understanding benchmarks, such as ARC-Easy~\cite{arc}, ARC-Challenge~\cite{arc}, Hellaswag~\cite{hellaswag}, Winogrande~\cite{winoGrande}, PIQA~\cite{piqa}, OpenbookQA~\cite{openbookqa}, and BoolQ~\cite{boolq}. Additionally, we provide validation perplexity figures on both WikiText2~\cite{wikitext2} and C4~\cite{c4} datasets.

We systematically measured GPU memory utilization and inference latency for both \text{LLaMA LLM} and \text{BitNet b1.58}. All experiments were run using the FasterTransformer\footnote{\url{https://github.com/NVIDIA/FasterTransformer}} library, which offers highly optimized GPU inference for large language models. The experiments also incorporated Ladder's~\cite{ladder} 2-bit kernel into \text{BitNet b1.58}. The primary latency metric reported is the time taken per generated output token, reflecting the dominant cost during inference.

Table~\ref{tab:ppl} presents a summary of model perplexities alongside the corresponding resource consumption for \text{BitNet b1.58} and \text{LLaMA LLM}. Notably, \text{BitNet b1.58} achieves comparable perplexity to full-precision \text{LLaMA LLM} at the 3B scale while providing a 2.71-fold speedup and a 3.55-fold reduction in GPU memory footprint. At 3.9B parameters, \text{BitNet b1.58} further exceeds the performance of the 3B \text{LLaMA LLM}, operating 2.4 times faster and consuming only about one-third of the memory.

In Table~\ref{tab:zero-shot}, we provide a comprehensive breakdown of zero-shot task accuracies. Evaluation procedures were aligned with the protocol implemented in the \emph{lm-evaluation-harness}\footnote{\url{https://github.com/EleutherAI/lm-evaluation-harness}} framework. The data indicate that as model scale increases, the difference in accuracy between \text{BitNet b1.58} and \text{LLaMA LLM} decreases. Crucially, from 3B parameters onward, \text{BitNet b1.58} achieves parity with full precision baselines. Echoing the perplexity findings, the task-level outcomes confirm that \text{BitNet b1.58} at 3.9B parameters surpasses \text{LLaMA LLM} 3B in both effectiveness and resource efficiency, underscoring that \text{BitNet b1.58} represents a clear Pareto improvement over leading large language model architectures.

\paragraph{Memory and Latency}

To explore scalability, we increased the model size to 7B, 13B, and 70B parameters, and assessed the associated computational cost. Figure~\ref{fig:latency-memory-energy} displays the patterns observed for both latency and memory usage, indicating that acceleration becomes more pronounced with larger models. Notably, \text{BitNet b1.58} 70B achieves a 4.1-fold speed advantage compared to the \text{LLaMA LLM} reference. This improvement stems from the fact that the computational burden of \emph{nn.Linear} layers escalates as model size grows. The pattern in memory utilization is analogous, since the embedding layer retains full precision and constitutes a decreasing share of overall memory consumption as models become larger. Both latency and memory statistics were gathered using a 2-bit kernel implementation, which suggests that further optimizations are possible to drive costs even lower.

\paragraph{Energy}

We conducted a comparative estimation of the energy usage resulting from arithmetic operations in both \text{BitNet b1.58} and \text{LLaMA LLM}. Our primary attention was on matrix multiplication computations, as they represent the dominant component of LLM computational expense. Figure~\ref{fig:energy} depicts the breakdown of energy usage, where the majority of the operations in \text{BitNet b1.58} consist of INT8 additions, in contrast to the combination of FP16 additions and multiplications found in \text{LLaMA LLM}. Leveraging the energy consumption models presented in~\cite{energycost, pokebnn}, we observe that \text{BitNet b1.58} delivers a 71.4-fold reduction in energy required for matrix multiplication arithmetic on 7nm process technology. Additionally, we provide end-to-end energy usage figures for processing sequences of 512 tokens. Our findings demonstrate that as model size increases, the energy efficiency of \text{BitNet b1.58} surpasses that of the FP16 \text{LLaMA LLM} baseline by greater margins. This is attributed to the fact that \emph{nn.Linear} operations comprise an increasing share of computational cost in larger models, while contributions from other layers diminish correspondingly.

\paragraph{Throughput}

We evaluate and contrast the throughput of \text{BitNet b1.58} and \text{LLaMA LLM} models, each comprising 70B parameters, deployed on a pair of 80GB A100 GPUs utilizing pipeline parallelism~\cite{gpipe}. This configuration enables the execution of the \text{LLaMA LLM} 70B model across the available hardware. Batch size was incrementally scaled up to the threshold permitted by GPU memory, while the input sequence length was fixed at 512 tokens. According to Table~\ref{tab:throughput}, the \text{BitNet b1.58} 70B model can process a batch size up to 11 times larger than that of the \text{LLaMA LLM}, yielding an increase in throughput by a factor of 8.9.

\textbf{\text{BitNet b1.58} introduces an alternative scaling paradigm connecting model capability and inference efficiency}. For illustrative purposes, equivalencies between model capacities for 1.58-bit and standard 16-bit representations, as evidenced in Figure~\ref{fig:latency-memory-energy} and \ref{fig:energy}, are summarized below.

\begin{itemize}
    \item A 13B parameter BitNet b1.58 outperforms a 3B FP16 LLM in terms of latency, memory demands, and energy efficiency.
    \item A 30B parameter BitNet b1.58 demonstrates lower latency, reduced memory usage, and decreased energy consumption compared to a 7B FP16 LLM.
    \item Similarly, a 70B BitNet b1.58 model offers better efficiency regarding latency, memory, and energy than a 13B FP16 LLM.
\end{itemize}

\paragraph{Training with 2T Tokens}

The scale of training tokens fundamentally impacts the effectiveness of large language models. To assess how well \text{BitNet b1.58} accommodates extensive token counts, we trained this model variant using 2T tokens, adhering to the data preparation methodology of StableLM-3B~\cite{StableLM-3B-4E1T}, which currently represents the leading 3B-parameter open-source model. Both systems were evaluated using a test suite that comprises Winogrande~\cite{winoGrande}, PIQA~\cite{piqa}, SciQ~\cite{sciq}, LAMBADA~\cite{lambada}, and ARC-easy~\cite{arc}. The zero-shot performance metrics are reported in Table~\ref{tab:2t}. For those benchmarks using both accuracy and normalized accuracy, we computed the mean across these two measures. StableLM 3b results for 2T tokens were directly obtained from its official technical documentation. Our analysis demonstrates that \text{BitNet b1.58} consistently outperforms on all considered benchmarks, highlighting robust generalization performance by 1.58-bit LLM architectures.

\section{Discussion and Future Work}

\textbf{1-bit Mixture-of-Experts (MoE) LLMs}

Mixture-of-Experts (MoE) architectures have established themselves as an efficient solution for scaling LLMs. Although MoE substantially decreases computational FLOPs, the substantial memory requirements and the cost of communication between processing units present obstacles to their widespread adoption. Employing 1.58-bit LLMs provides avenues to mitigate these problems. The memory requirements are lessened, which directly reduces the hardware needed to support MoE architectures. Furthermore, lower memory usage diminishes the network load involved in exchanging activations between chips. Ideally, this makes it feasible to fit entire models onto a single device, eliminating communication overhead altogether.

\textbf{Native Support of Long Sequence in LLMs}

Addressing long-context processing is increasingly essential for contemporary LLMs. Memory usage from maintaining KV caches constitutes a principal bottleneck for long sequence inference. \text{BitNet b1.58} offers an advance by decreasing activation storage from 16 bits to 8 bits, thus effectively doubling the attainable context window under fixed resources. Future extensions could focus on further, lossless compression—potentially achieving 4 bits per activation or even less for 1.58-bit LLMs—which remains an open avenue for continued investigation.

\textbf{LLMs on Edge and Mobile}

Deploying 1.58-bit LLMs offers transformative benefits for language modeling on resource-constrained edge and mobile hardware. These platforms face significant memory and computation limitations that impede large-scale LLM usage. By lowering both energy and memory demands, 1.58-bit LLMs enable deployment on devices previously unsuitable for advanced LLM applications, broadening their scope and functionality. This advancement paves the way for innovative on-device applications. Additionally, the 1.58-bit precision is well-matched to CPU-based devices, which dominate edge and mobile computing, thus allowing \text{BitNet b1.58} to execute efficiently and to further elevate device performance.

\textbf{New Hardware for 1-bit LLMs}

Efforts such as Groq\footnote{\url{https://groq.com/}} have showcased substantial promise in crafting specialized hardware (e.g., LPUs) optimized for LLM workloads. Building upon these advances, we propose the design and development of novel hardware and computing platforms that target the specific demands of 1-bit LLMs, taking advantage of the computation framework introduced by BitNet~\cite{bitnet}.